\begin{center}
[5.]
\end{center}
\begin{center}
\textsc{Gauss} an \textsc{Gerling}.\quad Göttingen, 5.~October 1821.
\end{center}

\par\noindent\hspace*{\parindent}Meinen herzlichen Glückwunsch zu dem von Ihnen übernommenen Geschäft.
Es freut mich, dass es in Ihre Hände kommt, da Sie es sich zur Pflicht machen werden, das ganze Triangelnetz mit aller möglichen Sorgfalt auszuführen.
Ich halte dies für etwas überaus wünschenswerthes, und beklage es immer, wenn man schon bei den Triangeln der ersten Ordnung geizig überlegt, durchaus nichts mehr an Genauigkeit anwenden zu wollen, als für den allerletzten Zweck unumgänglich nöthig ist.
Die \textls[120]{genaueste} Kenntniss der relativen Lagen der interessantesten Punkte eines Landes kann in vielfacher Beziehung nützlich sein, auch ganz abgesehen davon, dass eine Detailvermessung darauf am besten zu stützen ist.
Es wäre gewiss äusserst wichtig, wenn der grösste Theil von Europa vollständig mit Einem Netz überzogen wäre, und nach und nach werden wir dahin kommen; jeder Staat sollte es sich zur Ehre rechnen, seinen Antheil daran so gut zu liefern, dass er würdig sei, neben den besten zu stehen.

\par\noindent\hspace*{\parindent}Ich sollte glauben, dass Sie gut thun würden, den Meisner auch mit zu einem Dreieckspunkt zu wählen; Sie werden denselben \textls[120]{unmittelbar} mit wenigstens 5 meiner Dreieckspunkte in Verbindung setzen können, und wahrscheinlich mit ebenso viel \textsc{Müffling}schen dazu.
Ich werde gern dabei hülfreiche Hand bieten, und insofern es sich nur irgend einrichten lässt, die diesseitigen Beobachtungen dahin machen.
Die Heliotrope erleichtern solche Arbeiten ausserordentlich, und es bedarf auf dem Meisner weiter keiner Anlage, als dass ein steinernes Postament von circa \(1\frac{3}{4}\) Fuss Quadrat und 3 Fuss Höhe über der Erde gesetzt werde, um Heliotrop und Theodolithen darauf zu stellen, und dass vielleicht einige Bäume gefällt werden, falls ohne das die
\par\clearpage
\noindent\null\hfill{\small\textbf{ZUR HANNOVERSCHEN TRIANGULATION.}}\hfill\makebox[0pt][r]{\fontsize{13}{16}\selectfont \textbf{381}}\par\medskip
\noindent freie Aussicht nicht vollständig erreicht werden kann.\quad\mbox{.\ .\ .\ .\ .}\quad Von Hannover aus ist mir gar keine Communication Ihre Triangulirung betreffend zugegangen; ich habe daher den Köterberg, der sonst vortrefflich zum genauen Anschluss des Hercules gedient hätte, ausgeschlossen.
Inzwischen habe ich die Richtung der Linie Hohehagen-Hercules zu den andern dortigen sorgfältig gemessen.\\
\mbox{.\ .\ .\ .\ .\ .}
